% Figure 13 — truth recoverability across a four-stage degradation cascade.
% The curve is schematic; interventions are transverse cuts that stop propagation.
\begin{figure}[H]
\centering
\begin{tikzpicture}[pd figure,x=.94cm,y=.92cm]
  \draw[pd arrow] (.45,.35)--(15.15,.35);
  \draw[pd arrow] (.45,.35)--(.45,5.15);
  \node[pd axis label,anchor=west] at (13.15,.04) {cascade progression};
  \node[pd axis label,rotate=90,anchor=south] at (-.12,2.75)
    {recoverable truth};

  \draw[pd caution rule,line width=1.2pt]
    plot[smooth] coordinates {(.85,4.62)(4.15,4.00)(7.45,2.95)(10.75,1.55)(14.25,.72)};
  \foreach \x/\y/\n/\lab in {
    .85/4.62/1/context rot,
    4.15/4.00/2/lost in the middle,
    7.45/2.95/3/recursive collapse,
    10.75/1.55/4/over-flattening}{
    \node[pd caution datum] at (\x,\y) {};
    \node[pd panel title,anchor=north,align=center] at (\x,-.12) {\n\\\lab};
  }

  % Countermeasures cut the causal path at the stage they address.
  \foreach \x/\y/\short in {
    2.45/4.38/compact earlier,
    5.80/3.52/edge-place facts,
    9.05/2.28/from artifacts,
    12.55/1.08/zoom path}{
    \draw[pd focus rule,line width=1.5pt] (\x,\y-.42)--(\x,\y+.42);
    \node[pd direct label,text=hhteal,anchor=south,align=center] at (\x,\y+.52) {\short};
  }
  \draw[pd focus arrow] (9.05,2.80)--(9.05,4.65);
  \node[pd direct label,text=hhteal,anchor=south,align=center] at (9.50,4.72)
    {critical break\\reconstruct from evidence};
  \node[pd note,anchor=west] at (.85,.72) {schematic causal order; not a measured decay curve};
\end{tikzpicture}
\caption{The context-degradation cascade as loss of recoverable truth. Context
rot creates pressure to compact; middle loss removes paid-for facts; recursive
summarization fabricates; and the final digest replaces the artifact. A
countermeasure can interrupt each transition. The decisive cut is
compact-from-artifacts: rebuilding from durable, addressable evidence prevents
an earlier summary from becoming the next summary’s ground truth.}
\label{fig:cascade}
\end{figure}
